\begin{claim}{C.85}{Proved}{A parametric cutoff and all positive homotopy groups in the radial cyclic model}
Retain the exact coordinate $j$, original disk,
operator $q=jT_{(1,0)}j^{-1}$, and topological
group $\mathscr P$ of C.84.
There is a continuous map
\[
 \begin{gathered}
 \mathcal C:\mathscr P\times[1,\infty]
                 \longrightarrow\mathscr P,\qquad
 \mathcal C_R(F)=\mathcal C(F,R),\\
 \mathcal C_\infty(F)=F,\qquad
 \mathcal C_R(\id)=\id.
 \end{gathered}                                                   \tag{C.85.1}
\]
For every finite $R$, $\mathcal C_R(F)$ agrees
with $F$ on $|y|\leq R$ and is the identity
outside a finite horizontal band. The band
can be chosen uniformly for $F$ in any
compact family and $R$ in any compact
finite interval. Moreover,
$\mathcal C_R(F)\to F$ as $R\to\infty$
uniformly on compact families in the
original closed-disk map-and-inverse topology.
Consequently every based sphere family in
the actual relative centralizer
$\mathscr C_Q$ has a based filling there:
\[
                 \pi_n(\mathscr C_Q,1)=0
                         \qquad(n\geq1).                         \tag{C.85.2}
\]
Every original boundary point and every
original integer-translation label is
retained. This statement concerns the
specified radial cyclic action.
\end{claim}

\begin{proof}
\emph{The actual upper component and its marked conformal lift.}
Write $z=x+iy$ and put
\[
 \begin{gathered}
 E(z)=e^{2\pi iz},\qquad
 m=m_R(F)=\max_{0\leq x\leq1}|F(x,R)-(x,R)|,\\
 M=1+m,\qquad B=2(M+1)=2m+4,\qquad L=100B.
 \end{gathered}                                                   \tag{C.85.3}
\]
All these quantities refer to the actual
plane coordinates of C.84.
The proper line $\Gamma=F(\mathbb R\times\{R\})$
is periodic under $T_{(1,0)}$ and lies in
$R-m\leq y\leq R+m$.
Its upper complementary component is
$\Omega=F(\{y>R\})$. Indeed, the sublinear
displacement condition makes the image of
the positive end approach the positive end,
which identifies this component exactly.

The curve
$\gamma(x)=E(F(x,R))$, $x\in\mathbb R/\mathbb Z$,
is a Jordan curve. Equality of two values
means that their plane images differ by an
integer translation; injectivity of $F$
and its exact translation relation give
equality of the source labels modulo one.
Its winding about zero is one, since the
real part of $F(x+1,R)$ exceeds that of
$F(x,R)$ by one. Its bounded Jordan domain
$D$ therefore contains zero.
The preimage of $\gamma$ under $E$ is
exactly $\Gamma$. Points sufficiently high
above $\Gamma$ map into $D$, so
$E^{-1}(D\setminus\{0\})=\Omega$.

Let $f:\overline{\mathbb D}\to\overline D$
be the conformal Jordan-domain map normalized
by
\[
                  f(0)=0,\qquad f(1)=E(F(0,R)).                  \tag{C.85.4}
\]
Existence and uniqueness follow from the
Riemann map and its Jordan boundary
extension, followed by the unique rotation
of the source circle fixing the displayed
boundary value. The function
$g(w)=f(w)/w$, with $g(0)=f'(0)$,
is holomorphic and nonzero in the disk
and is continuous and nonzero on its
closure. Choose its continuous logarithm
by the exact equation
$\log g(1)=2\pi iF(0,R)$.
Define, for $\operatorname{Im}z\geq0$,
\[
 \begin{gathered}
 \phi(z)=z+\frac{1}{2\pi i}\log g(E(z)),\\
 d(z)=\phi(z)-z-iR.
 \end{gathered}                                                   \tag{C.85.5}
\]
The displayed logarithm exists because
$g$ is nonzero on the simply connected
closed disk; its holomorphic logarithm
in the interior extends continuously
to the boundary. Formula (C.85.5) fixes
the entire lift, including its integer
label: $\phi(0)=F(0,R)$ and
$\phi(z+1)=\phi(z)+1$.

It is a conformal bijection from the
upper half-plane onto $\Omega$.
For injectivity, equality of its values
first gives equality of $E(z)$ by the
injectivity of $f$, and then the period
relation removes the resulting integer.
For surjectivity, given $w\in\Omega$,
lift $f^{-1}(E(w))$ to the upper half-plane
and change that lift by the unique integer
making its $\phi$-image equal to $w$.
The same formulas give its continuous
boundary extension and inverse.

There is an increasing real homeomorphism
$\beta$ satisfying
\[
 \begin{gathered}
 \phi(x)=F(\beta(x),R),\qquad
 \beta(0)=0,\qquad \beta(x+1)=\beta(x)+1,\\
 \alpha=\beta^{-1},\qquad
 |\beta(x)-x|\leq1,\qquad |\alpha(x)-x|\leq1.
 \end{gathered}                                                   \tag{C.85.6}
\]
These statements follow by applying the
actual inverse of $F$ to the boundary
map. Its increasing orientation follows
from the period-one relation for a real
homeomorphism. The two bounds follow on
$[0,1]$ from the fixed endpoints and then
on all of $\mathbb R$ by periodicity.

The function $d$ in (C.85.5) is the
pullback of a holomorphic function on
the disk, continuous on its closure;
in particular the end $y=\infty$ is
its removable point $w=0$.
On the real boundary,
$|d(x)|\leq m+1=M$ by (C.85.6).
The maximum principle and the Cauchy
estimate on a radius-$y/2$ disk give
\[
             |d(z)|\leq B,\qquad
             |d'(z)|\leq\frac{2B}{y}\quad(y>0).                 \tag{C.85.7}
\]
Thus the bound includes the entire
original half-plane, without a new
boundary compactification assumption.

\emph{A cutoff which remains a homeomorphism.}
Fix the following $C^1$ function:
\[
 \chi(s)=
 \begin{cases}
  0,&s\leq1,\\
  3(s-1)^2-2(s-1)^3,&1\leq s\leq2,\\
  1,&s\geq2.
 \end{cases}
 \qquad 0\leq\chi\leq1,\quad |\chi'|\leq2.
                                                                    \tag{C.85.8}
\]
Put
\[
            \widetilde\phi(z)
             =z+iR+(1-\chi(y/L))d(z).                            \tag{C.85.9}
\]
For $0<y<L$ this is the original conformal
map $\phi$. For $y\geq L$, the operator
norm of the real derivative of its
perturbation of $z+iR$ is at most
\[
         |d'(z)|+\frac{2|d(z)|}{L}
                    \leq\frac{4B}{L}=\frac1{25}<1.              \tag{C.85.10}
\]
Hence its derivative is invertible there
and preserves orientation: the straight
segment from the identity derivative to
this derivative consists of invertible
matrices. The cutoff has matching first
derivatives at $L$ and $2L$.
It follows that $\widetilde\phi$ is a
local homeomorphism everywhere in the
open half-plane.

Its image stays in $\Omega$. Below $L$
this follows from $\widetilde\phi=\phi$;
above $L$ its height is at least
$R+L-B>R+m$, above the entire boundary
curve. It is proper as a map to $\Omega$.
Indeed, its distance from $z+iR$ is
bounded by $B$. An unbounded source
sequence therefore has unbounded image.
A bounded source sequence escaping
compact subsets of the open half-plane
has a subsequence approaching its real
boundary, whose image approaches
$\Gamma$ and cannot remain in a compact
subset of $\Omega$.

A proper local homeomorphism has an
open and closed image here and is a
covering onto that image. To recall
the latter assertion directly, a fiber
is compact and discrete, hence finite;
choose disjoint local inverse charts
at its points. Properness excludes
additional preimages over a sufficiently
small target neighborhood, giving an
evenly covered neighborhood.
Connectedness of $\Omega$ gives
surjectivity. A target of height
greater than $R+2L+B$ has exactly one
preimage: every preimage has $y>2L$,
where (C.85.9) is exactly $z+iR$.
The covering therefore has one sheet.
Its boundary extension is a
homeomorphism of the corresponding
closed half-planes as well. It is a
continuous bijection, and the same
bounded-distance estimate makes it
proper, including on the real boundary.
Its inverse is consequently continuous.

Set $\eta(y)=\min\{y,1\}$ for $y\geq0$ and
\[
 \begin{gathered}
 A(x,y)=\big((1-\eta(y))\alpha(x)+\eta(y)x,\ y\big),\\
 E^+_{F,R}(x,R+y)=\widetilde\phi(A(x,y)).
 \end{gathered}                                                   \tag{C.85.11}
\]
Each first-coordinate function of $A$
is increasing, proper, and period one.
Its inverse is continuous by the
two-sided monotone bracketing argument
of C.84. Thus $A$ is a homeomorphism
of the closed half-plane, commutes with
the original unit translation, and
has displacement at most one.
The inverse of $E^+_{F,R}$ is explicitly
the inverse of $\widetilde\phi$, followed
by $A^{-1}$ and by $(x,y)\mapsto(x,R+y)$.
It is a homeomorphism onto
$\overline\Omega$, and
\[
 \begin{gathered}
 E^+_{F,R}(x,R)=F(x,R),\\
 E^+_{F,R}(x,y)=(x,y)\quad(y\geq R+2L),\\
 |E^+_{F,R}(x,y)-(x,y)|\leq B+1=2m+5.
 \end{gathered}                                                   \tag{C.85.12}
\]
The identity above $R+2L$ uses $2L>1$,
so both cutoffs have already become
the identity there.

Let $J(x,y)=(x,-y)$. This exact reflection
commutes with $T_{(1,0)}$, and
$F\mapsto JFJ$ preserves $\mathscr P$
and its weighted metric. Define
$E^-_{F,R}=J E^+_{JFJ,R}J$ on $y\leq-R$.
The complete cutoff is
\[
 \mathcal C_R(F)(x,y)=
 \begin{cases}
   E^-_{F,R}(x,y),&y\leq-R,\\
   F(x,y),&-R\leq y\leq R,\\
   E^+_{F,R}(x,y),&y\geq R.
 \end{cases}                                                     \tag{C.85.13}
\]
The shared values on both boundary
lines agree. The images are respectively
the lower component, the middle strip
image under $F$, and the upper component
of the two actual image lines.
They cover the plane with disjoint
interiors. The inverses just described,
together with $F^{-1}$ on the middle
image, therefore paste to the actual
inverse of (C.85.13).
Every piece commutes with the original
unit translation. Outside a finite
band the map is the identity; on that
band its displacement is bounded by
periodicity. It thus belongs to
$\mathscr P$ by C.84. For $F=\id$,
the normalized conformal map is
$f(w)=e^{-2\pi R}w$, so $d=0$,
$\alpha=\id$, and every piece in
(C.85.13) is exactly the identity.

\emph{Finite-parameter continuity, with the chosen labels.}
Here is the precise conformal parameter
input. Rad\'o's original theorem
\cite[Section~1, p.~182; proof pp.~183--186]{C_Rado1923Variable}
states that normalized Riemann maps
with value zero and positive derivative
there converge uniformly on the entire
closed disk exactly when their enclosing
Jordan curves converge in Fr\'echet
distance. We use that theorem for the
specific Jordan curves $\gamma$ above.
The remaining normalization and topology
comparisons are as follows.

Take $F_i\to F$ in $\mathscr P$ and
finite $R_i\to R$. Their maps and
inverses converge uniformly on compact
plane sets by C.84. Hence
$E(F_i(x,R_i))\to E(F(x,R))$
uniformly for $0\leq x\leq1$; this
is an explicit Fr\'echet correspondence.
Rad\'o's theorem gives uniform
convergence of the positive-derivative
maps on the closed disk. Their boundary
preimages of $E(F_i(0,R_i))$ converge
as well: any convergent subsequence of
those circle points must have image
$E(F(0,R))$, whose limiting preimage
is unique. Rotating by these precise
circle points proves uniform convergence
$f_i\to f$ with normalization (C.85.4).

The functions $g_i=f_i/w$ also converge
uniformly on the closed disk. Indeed,
the difference quotient $(f_i-f)/w$
has a removable value at zero and its
maximum modulus is bounded by
$\|f_i-f\|_\infty$ on the unit circle.
Since $g$ is nonzero on the compact
disk, $g_i/g\to1$ uniformly. The small
logarithm of this ratio tends uniformly
to zero. It equals $\log g_i-\log g$
with our actual branches: this equality
holds at $w=1$ for large $i$ by
$F_i(0,R_i)\to F(0,R)$, and hence
everywhere by connectedness.
Consequently $d_i\to d$ uniformly on
the entire closed half-plane.

The quantities $m_i,B_i,L_i$ converge.
The functions $\chi(y/L_i)$ converge
uniformly for $y\geq0$ to $\chi(y/L)$:
they agree above one common finite
height, and convergence is uniform on
the remaining compact interval.
Thus $\widetilde\phi_i\to\widetilde\phi$
uniformly on the closed half-plane.
On $[0,1]$ the exact boundary inverses
\[
       \beta_i(x)=\operatorname{pr}_1
                    F_i^{-1}(\phi_i(x))                          \tag{C.85.14}
\]
converge uniformly to $\beta$.
Monotone bracketing and periodicity then
give uniform convergence of
$\alpha_i=\beta_i^{-1}$ to $\alpha$
on $\mathbb R$ in their displacement
coordinates. The maps $A_i$ converge
uniformly to $A$ on the closed
half-plane. Each fixed
$\widetilde\phi$ is uniformly continuous
there: its perturbation is periodic
in $x$, continuous on one compact
vertical strip modulo that period,
and zero above $2L$.
It follows that
$\widetilde\phi_i A_i\to\widetilde\phi A$
uniformly.

The two reflected versions have the
same properties. When the interfaces
$\pm R_i$ move, the formulas agree
exactly on each interface, so their
pasting is jointly continuous in the
parameters and the source point.
All cutoffs in this convergent sequence
are the identity outside one common
horizontal band. Uniform convergence
on its compact period rectangle gives
weighted forward convergence in C.84.
The radial estimate gives uniform
convergence of the original disk maps.
Their inverses converge uniformly too:
for homeomorphisms $a_i,a$ of the
compact disk,
\[
 \|a_i^{-1}-a^{-1}\|_\infty
       \leq\omega_{a^{-1}}
                     (\|a_i-a\|_\infty),                        \tag{C.85.15}
\]
where $\omega_{a^{-1}}$ is a modulus
of continuity of the fixed inverse.
C.84 now gives convergence in the
full actual group topology. This
proves finite-$R$ continuity, including
both inverse topology maps.

\emph{The infinite interface and compact families.}
For a compact $K\subset\mathscr P$ put
\[
 a_K(R)=\sup_{\substack{F\in K,\ 0\leq x\leq1\\|y|\geq R}}
                  \frac{|F(x,y)-(x,y)|}{|y|}.
                                                                    \tag{C.85.16}
\]
C.84 proves $a_K(R)\to0$.
The two boundary displacement maxima
at $\pm R$ are at most $Ra_K(R)$.
On the middle band (C.85.13) equals
$F$. On either tail (C.85.12) and its
reflection therefore give
\[
 \sup_{F\in K}\|\mathcal C_R(F)-F\|_*
          \leq 3a_K(R)+\frac{5}{1+R}
                   \longrightarrow0.                            \tag{C.85.17}
\]
The constants $c_F$ from (C.84.7) are
bounded uniformly for $F\in K$:
the common tail bound gives
$|Y_F(x,y)|\geq|y|/2$ off one band,
and the numerator there is bounded
on the remaining band. Formula
(C.84.8) thus turns (C.85.17) into
uniform forward convergence of the
original disk maps, uniformly over
$K$. The compact family of original
inverses has a common modulus of
continuity. Applying (C.85.15) with
that common modulus proves the same
uniform assertion for inverses.

This also proves joint continuity
at $R=\infty$. For a sequence
$F_i\to F$, $R_i\to\infty$, the set
$\{F,F_1,F_2,\ldots\}$ is compact,
so the estimate just proved gives
$\mathcal C_{R_i}(F_i)\to F$ in the
actual topology. Together with the
finite-parameter argument this proves
(C.85.1). In particular compact
quotient-support elements are dense,
with the displayed continuous cutoff
maps providing the parameter comparison.

\emph{The based sphere and its exact fixed-support filling.}
Let $a:(S^n,*)\to(\mathscr C_Q,1)$,
$n\geq1$, and write $F_\xi=j^{-1}a_\xi j$.
For $0<s\leq1$ use
$j\mathcal C_{1/s}(F_\xi)j^{-1}$,
and at $s=0$ use $a_\xi$.
Formula (C.85.1) proves that this is
a continuous based homotopy in the
original source group. Its last
family $a'_\xi=j\mathcal C_1(F_\xi)j^{-1}$
is supported, in the quotient
$N=\mathbb R^2/\langle T_{(1,0)}\rangle$,
in the interior of one fixed compact
annulus $N_0$. Such an annulus exists
because the two values of $L$ in
(C.85.3) are bounded on the compact
sphere family. Choose its two boundary
circles strictly beyond that common
support band.

Apply the exact compact-support
sphere filling of C.81 to the
descended family $b_\xi$ on $N_0$.
Its hypotheses hold here: the actual
extended plane translations form a
transitive boundary-fixed group
commuting with the original $Q$, as
in C.83; the quotient support is the
specified compact annulus with its
entire boundary fixed. In the notation
of (C.81.10), this gives
$B_{0,\xi}=1$, $B_{1,\xi}=b_\xi$,
$B_{r,*}=1$. Lift $B_{r,\xi}\pi$ from
the identity, where
$\pi:\mathbb R^2\to N$ is the original
covering. Lifting the inverse family
proves the inverse equations; uniqueness
of lifts proves commutation with every
original integer translation.
At $r=1$ the difference from
$\mathcal C_1(F_\xi)$ is an integer
deck label, locally constant on $S^n$
and zero at $*$. It is zero everywhere.
The whole lifted family and its inverse
fix the complement of the same band.
C.81's original-boundary extension,
or the bounded periodic displacement
estimate here, gives their joint
closed-disk continuity.

Writing this exact lifted disk family
as $H_{r,\xi}$, a complete based
null-homotopy is
\[
 K_{s,\xi}=
 \begin{cases}
   a_\xi,&s=0,\\
   j\mathcal C_{1/(2s)}(F_\xi)j^{-1},&0<s\leq\tfrac12,\\
   H_{2-2s,\xi},&\tfrac12\leq s\leq1.
 \end{cases}                                                     \tag{C.85.18}
\]
The two middle values agree with the
same original family $a'_\xi$.
The basepoint stays the identity,
$K_{0,\xi}=a_\xi$, and $K_{1,\xi}=1$.
This proves (C.85.2) in the original
topological group, with both maps
and inverses controlled throughout.

The conclusion includes all compact
based sphere families for this exact
radial cyclic action. The construction
uses its literal period-one coordinate,
its radial boundary topology, and the
sublinear comparison of C.84. The
general discrete-operator deformation
and general C.12 remain open.
\end{proof}
